\documentclass[11pt,a4paper]{article}
\usepackage[utf8]{inputenc}
\usepackage[french]{babel}
\usepackage{amsmath, amssymb, amsthm}
\usepackage{geometry}
\geometry{margin=1in}

\title{Analyse Structurale et Stratégie de Preuve Constructive de la Conjecture d'Erd\H{o}s-Gy\'{a}rf\'{a}s}
\author{Charles EDOU NZE\thanks{Charles EDOU NZE, chercheur indépendant}}
\date{}

\newtheorem{theorem}{Théorème}[section]
\newtheorem{lemma}[theorem]{Lemme}
\newtheorem{definition}[theorem]{Définition}
\newtheorem{conjecture}[theorem]{Conjecture}

\begin{document}
\maketitle

\begin{abstract}
Ce document présente une formulation axiomatique, une revue de la littérature et une architecture de preuve structurée pour la conjecture d'Erd\H{o}s-Gy\'{a}rf\'{a}s en théorie des graphes. La conjecture stipule que tout graphe dont le degré minimum est d'au moins 3 contient un cycle dont la longueur est une puissance de 2. Nous décomposons le problème en plusieurs lemmes intermédiaires et fournissons une architecture adaptée à la vérification formelle dans Lean 4.
\end{abstract}

\section{Introduction et Définitions Axiomatiques}
La conjecture d'Erd\H{o}s-Gy\'{a}rf\'{a}s est un problème ouvert fondamental en théorie des graphes extrémaux. Soit $G = (V, E)$ un graphe simple, fini et non orienté. Soit $\delta(G)$ le degré minimum des sommets dans $G$.

\begin{definition}[Graphe et Degré Minimum]
Soit $V$ un ensemble fini et $E \subseteq \binom{V}{2}$ l'ensemble des arêtes. Le degré minimum est défini par :
\[
\delta(G) = \min_{v \in V} \deg(v)
\]
où $\deg(v) = |\{u \in V \mid \{u, v\} \in E\}|$.
\end{definition}

\begin{conjecture}[Erd\H{o}s-Gy\'{a}rf\'{a}s, 1995]
Si $G = (V, E)$ est un graphe avec $\delta(G) \ge 3$, alors $G$ contient un cycle $C_k$ de longueur $k$ tel que $k = 2^m$ pour un certain entier $m \ge 1$.
\end{conjecture}

\section{Littérature Contextuelle et Analogies}
Le problème se situe à l'intersection de la théorie des graphes extrémaux et de la combinatoire additive. Il est structurellement similaire aux questions concernant l'existence de longueurs de cycles spécifiques dans les graphes denses (par exemple, les théorèmes de Bondy et Simonovits sur les cycles de longueur paire).

Nous dressons une analogie avec la preuve de l'existence de cycles pairs dans des graphes de densité suffisante, où l'on construit typiquement un chemin maximal et étudie le voisinage de ses extrémités. Ici, nous cherchons à affiner le processus de construction du chemin en exploitant la condition sur le degré pour forcer des cycles de longueurs spécifiques.

\section{Stratégie de Preuve et Lemmes}

Nous décomposons la conjecture en sous-problèmes structurés.

\begin{lemma}[Voisinage du Chemin Maximal]
\label{lem:longest_path}
Soit $P = v_1 v_2 \dots v_l$ un chemin maximal dans un graphe $G$ avec $\delta(G) \ge 3$. Alors tous les voisins de $v_1$ doivent se trouver sur $P$.
\end{lemma}

\begin{proof}
Supposons qu'il existe un voisin $u$ de $v_1$ tel que $u \notin V(P)$. Alors le chemin $u v_1 v_2 \dots v_l$ a une longueur $l+1$, contredisant l'hypothèse selon laquelle $P$ est un chemin maximal. Ainsi, $N(v_1) \subseteq V(P)$.
\end{proof}

\begin{lemma}[Formation de Cycles]
\label{lem:cycle_lengths}
Soit $P = v_1 v_2 \dots v_l$ un chemin maximal dans $G$ avec $\delta(G) \ge 3$. Soient $v_{i_1}, v_{i_2}, \dots, v_{i_k}$ les voisins de $v_1$ sur $P$, ordonnés tels que $2 = i_1 < i_2 < \dots < i_k \le l$. Puisque $\delta(G) \ge 3$, $k \ge 3$. Les arêtes $\{v_1, v_{i_j}\}$ créent des cycles $C^{(j)}$ de longueur $i_j$.
\end{lemma}

\begin{proof}
Le segment de chemin $v_1 \dots v_{i_j}$ a une longueur $i_j - 1$. L'ajout de l'arête $\{v_1, v_{i_j}\}$ complète un cycle de longueur $(i_j - 1) + 1 = i_j$.
\end{proof}

\section{Architecture de Formalisation Lean 4}
Pour autoformaliser cette preuve, nous esquissons les structures Lean 4.

\begin{verbatim}
import Mathlib.Combinatorics.SimpleGraph.Basic
import Mathlib.Combinatorics.SimpleGraph.Connectivity

variable {V : Type*} [Fintype V] [DecidableEq V]
variable (G : SimpleGraph V)

def min_degree_at_least (k : \mathbb{N}) : Prop :=
  \forall v : V, G.degree v \ge k

def has_cycle_power_of_two : Prop :=
  \exists (c : G.Walk v v), c.IsCycle \land \exists (m : \mathbb{N}), c.length = 2^m

theorem erdos_gyarfas (h : min_degree_at_least G 3) :
  has_cycle_power_of_two G := by
  sorry
\end{verbatim}

\section{Conclusion}
La décomposition structurelle fournie sert d'étape fondamentale vers la résolution de la conjecture d'Erd\H{o}s-Gy\'{a}rf\'{a}s, la préparant pour la vérification formelle.

\end{document}
